


    

        
    \ldtbox{
        \ldttitle{General trust and legitimacy}
        \ldtsubheading{To what extent:}
        \begin{itemize}[itemsep=1pt,topsep=0pt,leftmargin=8mm]
          \item[\arrowbox{↑↑}] Is trust and legitimacy important to the unilateral authority?
          \item[\arrowbox{↑↑}] Is the unilateral authority not trusted by other actors, regardless of the decisions it makes?
        \end{itemize}
                
        \vspace{0.2cm}

        \ldtsubheading{To what extent does the decision involve:}
        \begin{itemize}[itemsep=1pt,topsep=0pt,leftmargin=8mm]
          \item[\arrowbox{↑↓}] No clear expert consensus?
          \item[\arrowbox{↑}] Significant public interest concerns or externalities?
          \item[\arrowbox{↑↓}] Limited public impact?
          \item[\arrowbox{↑↓}] No clear expert consensus?
          \item[\arrowbox{↓}] A private technical or operational matter?
        \end{itemize}
    }

    \ldtbox{
        \ldttitle{External stakeholder legitimacy}
        \ldtsubheading{To what extent is there:}
        \begin{itemize}[itemsep=1pt,topsep=0pt,leftmargin=8mm]
            \item[\arrowbox{↑}] Powerful stakeholder groups, with conflicting perspectives?
            \item[\arrowbox{↑}] A bias for inaction given such conflict?
            \item[\arrowbox{↑}] An opportunity for decreased criticism or consequences due to ``process legitimacy''
        \end{itemize}
    }

    \ldtbox{
        \ldttitle{Internal stakeholder legitimacy}
        \ldtsubheading{To what extent is there:}

        \begin{itemize}[itemsep=1pt,topsep=0pt,leftmargin=8mm]
            \item[\arrowbox{↑}] Significant internal disagreement?
            \item[\arrowbox{↑}] A need to evolve organizational values or purpose with a strong mandate?
            \item[\arrowbox{↑}] Talent motivated by public benefit?
            \item[\arrowbox{↑}] A history of collaborative decision-making?
            \item[\arrowbox{↑↓}] Hierarchical structure
            \begin{itemize}[topsep=0pt,leftmargin=8mm]
                \item[\arrowbox{↑}] Can make it easier to delegate power to a democratic system
                \item[\arrowbox{↓}] Can correspond to reduced respect for such systems
            \end{itemize}
            \item[\arrowbox{↑↓}] Extensive collaboration between teams / departments?
            \begin{itemize}[topsep=0pt,leftmargin=8mm]
                \item[\arrowbox{↑}] Can correspond with flexible systems
                \item[\arrowbox{↓}] Can indicate complex interdependencies
            \end{itemize}
        \end{itemize}
    }

    \ldtbox{
        \ldttitle{Government risk}
        \ldtsubheading{To what extent is there:}

        \begin{itemize}[itemsep=1pt,topsep=0pt,leftmargin=8mm]
            \item[\arrowbox{↑}] Risk of backlash (e.g., antitrust, regulatory) if power is too concentrated?
            \item[\arrowbox{↑}] Pressure from politicians or autocrats to act in a way clearly against the wishes of the public?
            \item[\arrowbox{↑}] Limited oversight?
            \item[\arrowbox{↑↓}] High scrutiny environment?
            \begin{itemize}[topsep=0pt,leftmargin=8mm]
                \item[\arrowbox{↑}] Proactive democratization can accelerate preferential regulatory outcomes
                \item[\arrowbox{↓}] Lack of regulator experience with democratic processes or systems may limit influence
            \end{itemize}
            \item[\arrowbox{↑↓}] Complex regulatory landscape?
            \item[\arrowbox{↓}] Clear existing law?
        \end{itemize}
    }

    \ldtbox{
        \ldttitle{Collective intelligence}
        \ldtsubheading{To what extent does the decision:}

        \begin{itemize}[itemsep=1pt,topsep=0pt,leftmargin=8mm]
            \item[\arrowbox{↑}] Benefit from a broad diversity of perspectives?
            \item[\arrowbox{↑}] Benefit from locally distributed knowledge?
            \item[\arrowbox{↑}] Involve high levels of uncertainty?
            \item[\arrowbox{↓}] Involve complex interdependencies?
        \end{itemize}
    }

    \ldtbox{
        \ldttitle{Viability}
        \ldtsubheading{To what extent:}

        \begin{itemize}[itemsep=1pt,topsep=0pt,leftmargin=8mm]
            \item[\arrowbox{↑↓}] Are there political obstacles to delegating power?
            \begin{itemize}[topsep=0pt,leftmargin=8mm]
                \item[\arrowbox{↑}]Democratic legitimacy could resolve barriers
                \item[\arrowbox{↓}]Political obstacles are despite democratic norms
            \end{itemize}
        \item[\arrowbox{↓}] Are there organizational structure, legal, technical, or physical obstacles to further delegating or devolving power?
        \end{itemize}
    }

    \ldtbox{
        \ldttitle{Resources}
        \ldtsubheading{To what extent:}

        \begin{itemize}[itemsep=1pt,topsep=0pt,leftmargin=8mm]
            \item[\arrowbox{↑}] Is the level of resourcing made available for the democratic system commensurate with the importance of its decisions?
            \item[\arrowbox{↑}] Can resources be made available for recurring expenses (only impacts L4, L5)?
        \end{itemize}
    }

    \ldtbox{
        \ldttitle{Speed}
        \ldtsubheading{To what extent does the decision involve:}

        \begin{itemize}[itemsep=1pt,topsep=0pt,leftmargin=8mm]
            \item[\arrowbox{↑↓}] Time-critical responses?
            \begin{itemize}[topsep=0pt,leftmargin=8mm]
                \item[\arrowbox{↑}] Can establish processes suited to time constraints
                \item[\arrowbox{↑}] L4 processes can respond to regular time-critical scenarios in specified ways
                \item[\arrowbox{↓}] Rules out certain complex systems
            \end{itemize}
            \item[\arrowbox{↓}] Emergency responses?
        \end{itemize}
    }

    \ldtbox{
        \ldttitle{Adaptability}
        \ldtsubheading{To what extent do you anticipate:}

        \begin{itemize}[itemsep=1pt,topsep=0pt,leftmargin=8mm]
            \item[\arrowbox{↓}]Rapid changes to internal or external conditions that might impact the decision, relevance, or remit of a process?
            \begin{itemize}[topsep=0pt,leftmargin=8mm]
                \item[\arrowbox{↑}] This may be counteracted by more repeating or adaptive L4 or L5 systems
            \end{itemize}
        \end{itemize}
    }

    \ldtbox{
        \ldttitle{Novelty}
        \ldtsubheading{To what extent:}

        \begin{itemize}[itemsep=1pt,topsep=0pt,leftmargin=8mm]
            \item[\arrowbox{↑↓}] Is devolving such decisions novel?
            \begin{itemize}[topsep=0pt,leftmargin=8mm]
                \item[\arrowbox{↑}] Democratic first-mover might increase respect and newsworthiness
                \item[\arrowbox{↓}] Exposed to unknown risks
            \end{itemize}
        \end{itemize}
    }
\clearpage
